\subsection{Relative finite-arity lifting within a dimension range}

\begin{theorem}[C29: \status{Proved}; relative bar extension]
\label{sa08:C29}
Let $q=2r-1>0$ be odd, let $K=K(\Z,n)$, and let
$K\hookrightarrow D=\mathrm{DK}(D_{n+1}\Z)$ be the standard
sphere-to-disk inclusion. Use the positive bar closure $U_r$ of
Statement~\ref{sa08:D029}. Suppose a full strict simplicial PD-DGA
map $U_r(K)\to A$ is prescribed and, compatibly with it, a map
from $U_r(D)$ is prescribed through arity $m-1$, where $m\geq2$.
Give the parameter weight $w$. Assume that in weight $mw$ the
coefficient comparison is an isomorphism on $H^0$ and has zero
positive vertical cohomology through degree $(q+1)m$, and that
all horizontal rows in form degrees $0,\ldots,(q+1)m$ are acyclic.
If $n\leq2q-1$, the prescribed maps extend through arity $m$.

Thus every fixed finite arity range extends when its finite target
bounds hold. The assertion retains the entire prescribed boundary
and arbitrary compatible earlier choices. It does not assert one
uniform weight bound for all arities, or produce the initial map
$U_r(K)\to A$ from a map on its original disk.
\end{theorem}
\begin{proof}
Write $V_m(M)$ for the arity-$m$ part of the PD subalgebra
of $U_r(M)$ generated by smaller arities, and set
\[
 Q_m(M)=U_r(M)_m/V_m(M).
\]
The signed Lyndon basis of Statement~\ref{sa08:C24} makes its
terms free and places them in form degrees
$qm,\ldots,(q+1)m-1$. The coefficient term of degree $(q+1)m$
is already generated by arity one.

The Euler-derivation argument of Statement~\ref{sa08:C27}
applies without change to this regrading. The special derivation
homotopies have degree $-1$, preserve multidegree and arity, and
descend to $Q_m$. Each variable multiplicity therefore annihilates
homology; a component with multiplicity one is contractible.
Multiplication by $m$ is nullhomotopic. Consequently, writing
$F^j(M)=H^j(Q_m(M))$, every nonzero basis component of $F^j(M)$
uses at most $\lfloor m/2\rfloor$ variables. Moreover
\[
 F^j(M)=0\qquad(j\leq qm).
\]
Indeed, the first possible cohomology is a subgroup of the free
first term, but it is killed by $m$.

A basis element of $K(\Z,n)_a$ is labelled by an $n$-element
cut set. If $a>n\lfloor m/2\rfloor$, every surviving multidegree
component misses a common cut position and lies wholly in one
degeneracy image. The same subcomplex and its homology are obtained
from the split free submodule on those basis vectors. Hence
\[
 N_aF^j(K)=0\quad(a>n\lfloor m/2\rfloor),\qquad
 \operatorname{pd}_{s\mathrm{Ab}}F^j(K)
       \leq n\lfloor m/2\rfloor+1.
\]
The projective-dimension bound uses the identity-disk resolution
of a bounded normalized complex and the fact that every abelian
group has projective dimension at most one over $\Z$.
For $n\leq2q-1$ and $m\geq2$, this bound is at most $qm$.
Since nonzero $F^j$ have $j\geq qm+1$, their contributions to
derived Hom give
\begin{equation}
 H^{t\geq0}R\operatorname{Hom}_{s\mathrm{Ab}}
       (Q_m(K),\underline\Z)=0.
 \tag{C-R1}
\end{equation}
This follows by the finite vertical cohomology filtration, since
$\operatorname{Ext}^{j+t}(F^j(K),\underline\Z)=0$ for $t\geq0$.
The weaker bound $n\leq2q$ similarly gives vanishing for $t>0$,
recovering the corresponding general-degree absolute assertion.

The actual relative arity quotient is
\[
 R_m=Q_m(D)/Q_m(K).
\]
At each simplex, choose a basis of $D$ extending one of $K$.
Statement~\ref{sa08:C24} makes the indecomposable generators
for $K$ a basis subset of those for $D$. Thus the displayed
inclusion is injective and its quotient has free terms. The same
basis proves
\[
 U_r(K)_m\cap V_m(D)=V_m(K).
\]
The given full boundary and smaller-arity maps therefore agree
on this intersection and define a map on
$U_r(K)_m+V_m(D)$, with quotient $R_m$. This is the relative
problem, including the whole prescribed boundary.

Each term of $Q_m(D)$ is horizontally contractible: apply its
reduced functor to an additive simplicial contraction of $D$.
Using the derived-Hom identity from
Statement~\ref{sa08:C22},
\[
 R\operatorname{Hom}_{s\mathrm{Ab}}(Y,\underline\Z)
       \simeq R\operatorname{Hom}_{\Z}(NY,\Z),
\]
and the finite cochain filtration, its derived Hom into
$\underline\Z$ vanishes. The short exact sequence
$0\to Q_m(K)\to Q_m(D)\to R_m\to0$ consequently gives
\begin{equation}
 R\operatorname{Hom}_{s\mathrm{Ab}}(R_m,\underline\Z)
 \simeq
 R\operatorname{Hom}_{s\mathrm{Ab}}(Q_m(K),\underline\Z)[-1].
 \tag{C-R2}
\end{equation}
In particular its degree-one group vanishes by (C-R1).

The terms of $R_m$ have projective normalization. Since the
stated target rows are horizontally acyclic, extensions by those
rows split termwise in $s\mathrm{Ab}$. The strict cochain-map
extension obstruction is thus in
$H^1\operatorname{Tot}\operatorname{Hom}_{s\mathrm{Ab}}
(R_m,A^{mw,*})$. The ordinary-to-derived comparison from
Statement~\ref{sa08:C22}, applied to $R_m[-1]$ in degrees
$qm+1,\ldots,(q+1)m$, identifies it with the zero group
$H^1R\operatorname{Hom}(R_m,\underline\Z)$.
An actual simplicial cochain extension exists. Every product or
PD identity whose output has arity $m$ was already fixed on
the prescribed submodule, so this is the required strict partial
PD-DGA map. Repetition proves the finite-range assertion.
\end{proof}

The dimension boundary in this argument is substantive. For
$q=1$, $n=2$, and $m=2$, the absolute complex $Q_2(K)$ is
\[
 \operatorname{Sym}^2K\longrightarrow\Gamma^2K
\]
in form degrees two and three. Multiplication is injective on
each free module, with cokernel $K/2$; off-diagonal basis terms
map with coefficient one and diagonal terms with coefficient two.
It follows that
\[
 \begin{split}
 H^1R\operatorname{Hom}(R_2,\underline\Z)
 &\cong H^0R\operatorname{Hom}(Q_2(K),\underline\Z)\\
 &\cong \operatorname{Ext}_{s\mathrm{Ab}}^3
             (K(\mathbb F_2,2),\underline\Z)
 \cong \operatorname{Ext}_{\Z}^1(\mathbb F_2,\Z)
 \cong\mathbb F_2.
 \end{split}
\]
This is a nonzero possible obstruction group. It does not prove
that any prescribed admissible target and partial map realize its
nonzero element. In particular it is not a counterexample to C.3.
The initial boundary map into a common source whose row still
contains the cycle to be killed remains a separate requirement.

\begin{corollary}[C36: \status{Proved}; coefficient extension and local slope]
\label{sa08:C36}
Let $q=2r-1>0$ and $m\geq2$. The relative assertion of
Statement~\ref{sa08:C29}, with $K=K(k,n)$ and
$D=\mathrm{DK}(D_{n+1}k)$, holds over every commutative ring
$k$ when $n\leq2q-1$. Over a $\Z_{(p)}$-algebra it holds
in the larger range $n<pq$.

For the absolute problem, a prescribed strict map from $U_r(K)$
through arity $m-1$ extends through arity $m$ when $n\leq2q$
over arbitrary $k$, or when $n\leq pq$ over a
$\Z_{(p)}$-algebra. In both problems assume the finite target
conditions of C29: in weight $mw$, vertical coefficient comparison
through degree $(q+1)m$ and horizontal acyclicity in form degrees
$0,\ldots,(q+1)m$.

No flatness of $k$ over $\Z$ is required. The relative assertion
retains the prescribed full boundary map and arbitrary compatible
smaller-arity choices. These are finite-arity statements: they neither
provide a uniform weight for all arities nor construct the initial
full boundary map. At an individual arity $m$ prime to $p$, the
same extension conclusions over a $\Z_{(p)}$-algebra hold without
a restriction on $n$, under the same finite target conditions.
\end{corollary}
\begin{proof}
Let $R=\Z$ or $\Z_{(p)}$, and let $Q_R$ denote the bounded
universal arity quotient $Q_m(K(R,n))$ used above. Its terms are
degreewise free over $R$. Evaluation of a projective resolution
in $s\operatorname{Mod}_R$ is an $R$-projective resolution:
the standard projective generators evaluate to free $R$-modules,
and evaluation preserves exact sequences and direct summands.
Consequently the degreewise flat terms of $Q_R$ have no higher
left-derived extension-of-scalars groups. The finite cochain
filtration gives
\[
 k\otimes_R Q_R\simeq k\otimes_R^L Q_R.
\]
Free-module base change for divided powers and the bar construction,
together with the basis-subset quotient of C24, identifies the left
side with $Q_k$. The relative quotient $R_m$ has the same
degreewise freeness and base-change properties. Restriction of
scalars is exact, so its left adjoint preserves projectives and
derived adjunction gives
\[
 R\operatorname{Hom}_{s\operatorname{Mod}_k}(Q_k,\underline k)
 \simeq
 R\operatorname{Hom}_{s\operatorname{Mod}_R}(Q_R,\underline k),
\]
and likewise for $R_m$. The PID projective-dimension estimates
in C29 vanish against every target $R$-module, including $k$.
They therefore give the asserted integral absolute and relative
ranges over arbitrary $k$.

For $R=\Z_{(p)}$, the multiplicity Euler homotopy contracts a
multidegree component whenever one occurring multiplicity is not
divisible by $p$. Each surviving component thus uses at most
$\lfloor m/p\rfloor$ variables. For $F^j=H^j(Q_R)$ the same
cut-set argument gives
\[
 \operatorname{pd}_{s\operatorname{Mod}_R}F^j
 \leq n\lfloor m/p\rfloor+1,
 \qquad F^j=0\quad(j\leq qm).
\]
The second assertion uses the torsionfree first term over the PID
$R$ and the total Euler homotopy, exactly as in C29. If $n\leq pq$,
the projective-dimension bound is at most $qm+1$, whereas nonzero
$F^j$ start in degree $qm+1$. It follows that
$H^{t}R\operatorname{Hom}(Q_R,\underline k)=0$ for $t>0$.
If $n<pq$ and $p\mid m$, the bound is at most $qm$, so the
same vanishing holds also for $t=0$. When $p\nmid m$, total
Euler contracts $Q_R$ since $m$ is a unit, without a dimension
restriction.

The relative disk quotient has derived Hom equal to that of the
sphere quotient shifted by $-1$, as in (C-R2); its degree-one
obstruction therefore vanishes in the stated relative range.
Derived base change transfers all these vanishings to $k$.
The quotient terms have projective normalized complexes over $k$,
since normalization is a direct summand of each degreewise free
term. The finite target conditions now apply the ordinary-to-derived
Hom comparison of C22, giving actual strict cochain extensions.
The smaller-arity product and divided-power identities are already
prescribed on the same submodule as in C29, so the extensions are
strict partial PD-DGA maps. This proves exactly the finite absolute
and full-boundary relative claims.
\end{proof}
